% Separate subsection: direct-definition auxiliary benchmark.
% The user selected this direct auxiliary backend after its benchmark.
\subsection{Direct fermion sewing: definition and level-ten benchmark}
\label{direct-fermion-definition}

The selected algorithm evaluates the auxiliary factor directly in its
fermion basis, while retaining central-charge recursion for both
Virasoro factors in the enlarged block.  This subsection derives the
direct auxiliary construction and records its measured cost for the
inserted genus-two graph.  The timing below concerns the auxiliary
factor alone and performs no comparison with another block; validation
of the integrated physical calculation is reported separately.

On the NS edge use the states
$\boldsymbol\Psi_{-\mathsf A}\mathbf1$, where $\mathsf A$ is a
strictly descending list of positive half-integers.  On each Ramond
edge use $\boldsymbol\Psi_{-\mathsf B}u^{\mathsf b}$, where
$\mathsf B$ is a strictly descending list of positive integers and
$\mathsf b=0,1$.  The symbol $|\mathsf B|$ denotes the sum of the
modes, whereas $\#\mathsf B$ counts them.  We retain the conventions
\begin{equation}
 \{\psi_r,\psi_s\}=\delta_{r+s,0},\qquad
 \psi_0u^{\mathsf b}=\frac{u^{1-\mathsf b}}{\sqrt2},\qquad
 \psi_r^\dagger=-\psi_{-r},\qquad
 \langle u^{\mathsf b}|u^{\mathsf b}\rangle=(-1)^{\mathsf b}.
\end{equation}
The Fock metric is diagonal.  Its diagonal entries are
$(-1)^{\#\mathsf A}$ in the NS sector and
$(-1)^{\#\mathsf B+\mathsf b}$ in the Ramond sector.

The incoming state on the additional sphere is at $z_3=0$ and carries
$\hat q_{2,1}^{|\mathsf B|}$; the outgoing state is at infinity and
carries $\hat q_{2,2}^{|\mathsf B'|}$.  Write their common parity and
the parities on the other two edges as
\begin{equation}
 \delta_1=\#\mathsf A,\qquad
 \delta_2=\#\mathsf B+\mathsf b
          =\#\mathsf B'+\mathsf b',\qquad
 \delta_3=\#\mathsf C+\mathsf c\pmod2.
\end{equation}
Both outer forms vanish unless
$\delta_1+\delta_2+\delta_3=0\pmod2$.

\paragraph{The complete middle insertion.}
At $z_3=1$ every term in
$\psi(z_3)=\sum_{m\in\mathbb Z}\psi_mz_3^{-m-1/2}$ has unit
coordinate factor.  Using
\begin{equation}
 \Theta\boldsymbol\Psi_{-\mathsf B}u^{\mathsf b}
 =(-1)^{\#\mathsf B+\mathsf b+1}
  \boldsymbol\Psi_{-\mathsf B}u^{1-\mathsf b},
\end{equation}
the zero-mode action is
\begin{equation}
 Q\Theta\psi_0\boldsymbol\Psi_{-\mathsf B}u^{\mathsf b}
 =\frac Q{\sqrt2}(-1)^{\mathsf b}
  \boldsymbol\Psi_{-\mathsf B}u^{\mathsf b}.
\end{equation}
For $k>0$, creation by $\psi_{-k}$, when $k\notin\mathsf B$, and
annihilation by $\psi_k$, when $k\in\mathsf B$, both give the sign
$(-1)^{\#\{m\in\mathsf B:m>k\}}$ in the descending basis.
The resulting word has one more or one fewer fermion.  The subsequent
$\Theta$ action therefore contributes $(-1)^{\delta_2}$ and changes
the ground label to $1-\mathsf b$.  Consequently the action
coefficients of the full insertion are
\begin{equation}
 [Q\Theta\psi(1)]_{(\mathsf B',\mathsf b'),(\mathsf B,\mathsf b)}
 =\begin{cases}
  \dfrac Q{\sqrt2}(-1)^{\mathsf b},
       &\mathsf B'=\mathsf B,\quad\mathsf b'=\mathsf b,\\[3pt]
  Q(-1)^{\#\{m\in\mathsf B:m>k\}+\delta_2},
       &\mathsf B'=\mathsf B\mathbin\triangle\{k\},
          \quad\mathsf b'=1-\mathsf b,\quad k>0,\\[3pt]
  0,&\text{otherwise}.
 \end{cases}
 \label{direct-fermion-middle}
\end{equation}
Here $\triangle\{k\}$ toggles the occupancy of $k$.  For each incoming
state, each nonzero mode connects to at most one outgoing Fock state.
Both transitions
preserve $\delta_2$, although their two propagation levels need not
agree.  Reversing a toggle leaves its sign unchanged: the number of
occupied modes larger than $k$ and the total parity are unchanged.

\paragraph{The state sum and its signs.}
In the following formula the subscript on
$[Q\Theta\psi(1)]$ means an action coefficient, as in
\eqref{direct-fermion-middle}, rather than a BPZ matrix element.
The latter contains the additional outgoing norm
$(-1)^{\delta_2}$.  The four inverse Fock norms multiply to
$(-1)^{\delta_1+2\delta_2+\delta_3}$, so including this outgoing
norm leaves $(-1)^{\delta_1+\delta_2+\delta_3}=1$ on every allowed
summand.  The NS contour sign and the original theta sign remain.
With the two outer forms in the order $(\NS,\R,\R)$, the complete
direct definition is therefore
\begin{align}
 &\mathbb F_{\mathsf F}[\Theta v_{1/2}]
 (q_1,\hat q_{2,1},\hat q_{2,2},q_3)\nonumber\\
 &=\sum_{\mathsf A,\mathsf B,\mathsf b,
                \mathsf B',\mathsf b',\mathsf C,\mathsf c}
 q_1^{|\mathsf A|}\hat q_{2,1}^{|\mathsf B|}
 \hat q_{2,2}^{|\mathsf B'|}q_3^{|\mathsf C|}
 \eta_1^{\delta_1}\eta_2^{\delta_2}\eta_3^{\delta_3}
 (-1)^{\delta_1+\delta_1\delta_2+\delta_1\delta_3+\delta_2\delta_3}
 \nonumber\\
 &\quad\times
 \rho_{\mathsf F}\!\left(
  \boldsymbol\Psi_{-\mathsf A}\mathbf1,
  \boldsymbol\Psi_{-\mathsf B}u^{\mathsf b},
  \boldsymbol\Psi_{-\mathsf C}u^{\mathsf c}\right)
 \rho_{\mathsf F}\!\left(
  \boldsymbol\Psi_{-\mathsf A}\mathbf1,
  \boldsymbol\Psi_{-\mathsf B'}u^{\mathsf b'},
  \boldsymbol\Psi_{-\mathsf C}u^{\mathsf c}\right)
 [Q\Theta\psi(1)]_{(\mathsf B',\mathsf b'),(\mathsf B,\mathsf b)}.
 \label{direct-fermion-sum}
\end{align}
The ground forms are
$\rho_{\mathsf F}(\mathbf1,u^0,u^0)=1$ and
$\rho_{\mathsf F}(\mathbf1,u^1,u^1)=i$.  They give the constant
$Q(1-\eta_2\eta_3)/\sqrt2$ directly.  In particular the second
split edge does not introduce an independent parity sign.

The outer forms are computed by moving one fermion contour at a time
between the punctures at infinity, one and zero.  For the three possible
target slots, the useful test sections are
\begin{equation}
 z^{k-1/2}(z-1)^{-1/2},\qquad
 z^{-1/2}(z-1)^{-n-1/2},\qquad
 z^{-n-1/2}(z-1)^{-1/2}.
\end{equation}
Their expansions give half-integer binomial coefficients and the
frame phases $(i,1,-i)$, together with the usual graded contour signs.
For the first section the target NS creation mode is $k-1/2$; in the
other two sections the target Ramond creation mode is $n$.  Remove
the largest occupied mode in the first nonempty slot and solve the
Ward relation for the original form.  Every remaining term has lower
total nonzero-mode level, so this procedure terminates at the displayed
ground forms.  Each local expansion is stopped only after its modes
annihilate every state on that leg.  The resulting bound is determined
by the occupied modes, rather than a fixed expansion order.
Computed three-point forms are stored and reused throughout the sum.

\paragraph{Exact rational arithmetic.}
No irrational arithmetic is needed after extracting $Q/\sqrt2$ from
the final block.  To see this, temporarily use $\sqrt2u^1$ in place
of $u^1$.  In this basis the two zero-mode coefficients are $1/2$
and $1$.  Strip the known phase $i^{\delta_2}$ from the outer form
and denote the stored value by $\rho_{\mathsf F}^{\mathrm{rat}}$.
Thus, with the same fermion words as above,
\begin{equation}
 \rho_{\mathsf F}
 =i^{\delta_2}2^{-(\mathsf b+\mathsf c)/2}
   \rho_{\mathsf F}^{\mathrm{rat}}.
\end{equation}
The stored ground values are $1$ and $2$.  The mode Ward relations
then involve only rational binomial coefficients and signs, so every
stored value is rational.  The two outer phases in a sewn term
multiply to $(-1)^{\delta_2}$.

For a diagonal transition, combining these phases, the factor
$(-1)^{\mathsf b}$ from $\Theta\psi_0$, and the two ground rescalings
gives the following normalized summand, apart from the common
propagation monomial and sewing sign already displayed in
\eqref{direct-fermion-sum}:
\begin{equation}
 \frac{(-1)^{\#\mathsf B}}{2^{\mathsf b+\mathsf c}}
  \left(\rho_{\mathsf F}^{\mathrm{rat}}
   (\mathsf A,\mathsf B,\mathsf C)\right)^2.
\end{equation}
For a toggle, $\mathsf b+\mathsf b'=1$.  The square root of two in
$[Q\Theta\psi(1)]/(Q/\sqrt2)$ cancels the remaining half-integer
ground-rescaling power, while the two factors
$(-1)^{\delta_2}$ cancel.  Its normalized summand is
\begin{equation}
 \frac{(-1)^{\#\{m\in\mathsf B:m>k\}}}{2^{\mathsf c}}
 \rho_{\mathsf F}^{\mathrm{rat}}(\mathsf A,\mathsf B,\mathsf C)
 \rho_{\mathsf F}^{\mathrm{rat}}(\mathsf A,\mathsf B',\mathsf C).
\end{equation}
The ground labels in these abbreviated arguments remain those in
\eqref{direct-fermion-sum}.  The temporary basis rescaling has
therefore canceled from the physical normalization.  The implementation
uses exact FLINT rational arithmetic and restores $Q/\sqrt2$ only
when a numerical value of the full block is requested.

\paragraph{Finite truncation and measured cost.}
For the balanced split
$\hat q_{2,1}=u\sqrt{q_2}$,
$\hat q_{2,2}=u^{-1}\sqrt{q_2}$, total level $N$ means
\begin{equation}
 |\mathsf A|+
 \frac{|\mathsf B|+|\mathsf B'|}{2}+|\mathsf C|\leq N.
\end{equation}
The stored exponent tuple is
$(2|\mathsf A|,|\mathsf B|,|\mathsf B'|,2|\mathsf C|)$;
its entries sum to at most $2N$.  Unequal split levels retain the
power $u^{|\mathsf B|-|\mathsf B'|}$.  At $N=10$ a single split
edge can reach level twenty.  To enumerate each nonzero-mode pair
once, start from its lower-level state and add the unique toggled
mode $k$.  The reverse removal supplies the transposed pair of
propagation powers with the same coefficient.  The diagonal part
is added separately.

A fresh-process run on September 10, 2026, using Python 3.14.3 on
macOS arm64, computed the complete normalized auxiliary series through
$N=10$ in $0.186281209\,\mathrm{s}$.  This construction time includes
the provider, Fock bases, Ward cache, all mode transitions, rational
normalization and coefficient accumulation.  Elapsed process time
through imports and the initial result save was
$0.290214667\,\mathrm{s}$; peak resident memory was
$46.765625\,\mathrm{MiB}$.  The saved result contains $4793$
nonzero exponent tuples and $9586$ nonzero rational parity
coefficients.  The calculation stored $7508$ three-point forms and
summed $20156$ nonzero state contractions.  Its coefficients and
timing metadata are recorded in
\path{Code/theta_fermion_ccy/results/direct_fermion_level10.json}.
The final small rewrite adding timing metadata is excluded from the
process-through-save measurement.  These measurements establish the
cost of this direct auxiliary calculation; they do not constitute a
numerical validation or a timing of the complete physical block.

Equation \eqref{direct-fermion-sum} gives the full auxiliary factor.
Its normalization relative to the reduced central-charge recursion is
explained in \S\ref{ccy-and-truncation}.

\paragraph{Other genera and plumbing channels.}
The ingredients in \eqref{direct-fermion-sum} are local: a fermion
Fock basis on each edge, its diagonal pairing, a three-point form on
each sphere, and the sparse action of any inserted fermion.  They can
be sewn on any pants graph with compatible NS and Ramond labels.
NS--NS--NS vertices use the corresponding fermion Ward forms;
NS--R--R vertices use the forms described above.  Spin transports
and graded reordering signs are taken from the chosen graph and
local coordinates.  The particular quadratic theta sign in
\eqref{direct-fermion-sum} belongs to the present channel and must
not be imposed unchanged on another graph.  A finite plumbing-level
cutoff again leaves a finite sum, and the local three-point data can
be reused.  This is an arbitrary-genus sewing construction, although
the current assembler and the reported runtime concern only the
specified four-edge genus-two graph.  Its computational cost can
grow with the genus and the graph.  This direct auxiliary evaluation
does not use an Ising Verma-module limit or an Ising central-charge
recursion.  The two ordinary Virasoro blocks in the enlarged numerator
continue to use the prescribed central-charge recursion.
